\documentclass[11pt,a4paper]{article}
\usepackage[utf8]{inputenc}
\usepackage[T1]{fontenc}
\usepackage[spanish,es-nodecimaldot]{babel}
\usepackage{geometry}
\geometry{top=2.5cm, bottom=2.5cm, left=2.5cm, right=2.5cm}
\usepackage{graphicx}
\usepackage{amsmath}
\usepackage{amssymb}
\usepackage{booktabs}
\usepackage{hyperref}
\usepackage{xcolor}
\usepackage{caption}
\usepackage{float}
\usepackage{fancyhdr}
\usepackage{titlesec}
\usepackage{listings}

% Definición de Colores USS
\definecolor{USSBlue}{RGB}{0,32,91}
\definecolor{USSGold}{RGB}{212,175,55}
\definecolor{DarkGray}{RGB}{80,80,80}

% Configuración de Hipervínculos
\hypersetup{
    colorlinks=true,
    linkcolor=USSBlue,
    citecolor=USSBlue,
    urlcolor=USSGold,
    pdftitle={Informe Final - Taller de Aptitudes Lógicas y Matemáticas},
    pdfauthor={Moisés, Fernando, Juan Pablo}
}

% Estilo de Páginas (Cabecera y Pie)
\pagestyle{fancy}
\fancyhf{}
\fancyhead[L]{\textcolor{USSBlue}{\textbf{Taller de Aptitudes Lógicas y Matemáticas}}}
\fancyhead[R]{\textcolor{USSBlue}{Informe Final - Hito Cami}}
\fancyfoot[C]{\thepage}
\fancyfoot[R]{\textcolor{DarkGray}{Universidad San Sebastián}}
\renewcommand{\headrulewidth}{0.8pt}
\renewcommand{\footrulewidth}{0.4pt}

% Configuración de Títulos de Sección
\titleformat{\section}
  {\color{USSBlue}\normalfont\Large\bfseries}
  {\thesection}{1em}{}
\titleformat{\subsection}
  {\color{USSBlue}\normalfont\large\bfseries}
  {\thesubsection}{1em}{}

% Configuración para código fuente Python en el anexo
\lstset{
    language=Python,
    basicstyle=\ttfamily\small,
    keywordstyle=\color{USSBlue}\bfseries,
    stringstyle=\color{USSGold},
    commentstyle=\color{gray}\itshape,
    numbers=left,
    numberstyle=\tiny\color{gray},
    stepnumber=1,
    numbersep=8pt,
    backgroundcolor=\color{lightgray!15},
    showspaces=false,
    showstringspaces=false,
    frame=single,
    rulecolor=\color{USSBlue!30},
    breaklines=true,
    breakatwhitespace=true,
    tabsize=4
}

\begin{document}

% ==========================================
% PORTADA
% ==========================================
\begin{titlepage}
    \centering
    
\includegraphics[width=0.35\textwidth]{logo_uss.png}\\[1cm]
    
    \textcolor{USSBlue}{\rule{\textwidth}{1.5pt}}\\[0.4cm]
    {\Huge \bfseries \textcolor{USSBlue}{Diseño, Modelamiento Cinemático y Análisis de Integridad Estructural de un Brazo Robótico Antropomórfico de 3 GDL}}\\[0.2cm]
    {\Large \textcolor{USSGold}{\textbf{Informe Técnico Final - Desafío Matemático e Ingenieril}}}\\[0.4cm]
    \textcolor{USSBlue}{\rule{\textwidth}{1.5pt}}\\[1.5cm]
    
    \textbf{Asignatura:} Taller de Aptitudes Lógicas y Matemáticas (TALM)\\
    \textbf{Institución:} Universidad San Sebastián\\
    \textbf{Facultad:} Facultad de Ingeniería y Tecnología\\
    \textbf{Carrera:} Ingeniería Civil\\[1.5cm]
    
    \begin{minipage}{0.45\textwidth}
        \begin{flushleft} \large
            \textbf{Estudiantes:}\\
            Moisés\\
            Fernando\\
            Juan Pablo
        \end{flushleft}
    \end{minipage}
    ~
    \begin{minipage}{0.45\textwidth}
        \begin{flushright} \large
            \textbf{Profesor(a):}\\
            Cami\\
            \textbf{Fecha de Entrega:}\\
            Junio, 2026
        \end{flushright}
    \end{minipage}\\[2.5cm]
    
    \vfill
    {\textcolor{DarkGray}{\small Concepción, Chile}}
\end{titlepage}

\newpage
\tableofcontents
\newpage

% ==========================================
% SECCIÓN 1: INTRODUCCIÓN
% ==========================================
% [VALIDACIÓN LÓGICA]: Esta sección cumple con el criterio de "Comunicación Matemática" (CM) al introducir de manera clara e hilada el contexto físico del brazo robótico de 3 GDL, relacionándolo con la anatomía humana y delineando los objetivos generales y específicos.
\section{Introducción}
El desarrollo de sistemas mecánicos complejos y la robótica industrial moderna exigen metodologías de diseño rigurosas donde convergen la cinemática de sólidos, la dinámica clásica y la resistencia de materiales. En este contexto, el modelamiento de manipuladores robóticos antropomórficos representa un campo de estudio fundamental, dado que imitan la movilidad y la versatilidad de las extremidades humanas. 

Un brazo robótico antropomórfico de tres grados de libertad (GDL) se basa estructuralmente en la analogía de la extremidad superior del ser humano. Los tres GDL fundamentales corresponden a los siguientes movimientos:
\begin{enumerate}
    \item \textbf{Cintura (Rotación de la Base / Yaw):} Corresponde al movimiento rotacional alrededor del eje vertical $z$. Permite orientar el brazo en el espacio de trabajo acimutal.
    \item \textbf{Hombro (Pitch de la Articulación Intermedia):} Movimiento de elevación alrededor de un eje horizontal. Permite el control del alcance radial y de la altura.
    \item \textbf{Codo (Pitch Relativo de la Articulación Distal):} Movimiento rotacional del segundo eslabón respecto al primero, en un plano coplanar al del hombro, simulando la articulación del codo humano.
\end{enumerate}

El análisis de la \textit{cinemática directa} consiste en determinar la posición espacial del efector final (mano robótica o carga de pago) a partir de los valores angulares de las articulaciones ($\theta_1, \theta_2, \theta_3$). Este mapeo del espacio articular al espacio operacional se modela matemáticamente a través de proyecciones trigonométricas o transformaciones homogéneas de Denavit-Hartenberg.

El propósito general de este informe es diseñar, validar y evaluar estructuralmente un brazo robótico antropomórfico sometido a condiciones extremas de operación, asegurando su integridad física ante esfuerzos axiales y cortantes inducidos por rotación a altas velocidades, además de analizar los efectos de la dilatación térmica del sistema.

\subsection{Objetivos}
\textbf{Objetivo General:}
\begin{itemize}
    \item Diseñar y validar analíticamente un brazo robótico antropomórfico de 3 GDL con un alcance máximo de 20~cm, determinando su viabilidad mecánica ante una velocidad angular límite de rotación equivalente a 140~km/h en el efector final.
\end{itemize}

\textbf{Objetivos Específicos:}
\begin{itemize}
    \item Derivar las ecuaciones analíticas de la cinemática directa del manipulador tridimensional mediante la analogía de la sombra radial.
    \item Determinar la distribución de tensiones centrífugas en la base y el codo del manipulador bajo el escenario crítico de rotación a 140~km/h.
    \item Evaluar el comportamiento mecánico de los eslabones de Aluminio 6061-T6 y el pasador de la articulación del codo de Acero Inoxidable 316, identificando posibles puntos de falla crítica.
    \item Calcular la velocidad tangencial máxima permisible para garantizar que el pasador trabaje dentro de sus límites elásticos de cortante admisible.
    \item Modelar el impacto de la dilatación térmica lineal en la longitud de los eslabones y en el juego mecánico del pasador del codo ante un gradiente de temperatura de $\Delta T = 100^\circ$C.
\end{itemize}

% ==========================================
% SECCIÓN 2: DEFINICIÓN DEL PROBLEMA
% ==========================================
% [VALIDACIÓN LÓGICA]: Esta sección satisface el criterio de "Definición del Problema" (RP) y "Variables y Parámetros" (MM) al identificar y clasificar rigurosamente todas las variables físicas, parámetros de diseño y exigencias críticas del sistema.
\section{Definición del problema}
El sistema mecánico bajo análisis consiste en un brazo robótico antropomórfico tridimensional compuesto por dos eslabones móviles principales conectados en serie. Los parámetros geométricos e inerciales del diseño base se definen a continuación:
\begin{itemize}
    \item \textbf{Longitudes de los eslabones:} $L_1 = 10\text{ cm} = 0.1\text{ m}$ y $L_2 = 10\text{ cm} = 0.1\text{ m}$. El alcance máximo de extensión radial es $R_{\text{max}} = L_1 + L_2 = 20\text{ cm} = 0.2\text{ m}$.
    \item \textbf{Distribución de masas:}
    \begin{itemize}
        \item Cada eslabón tiene una masa distribuida uniformemente de $m_1 = m_2 = 0.5\text{ kg}$. Esto equivale a una densidad lineal de masa de $\rho = 5.0\text{ kg/m}$.
        \item El efector final soporta una carga de pago concentrada de masa $M = 0.5\text{ kg}$.
    \end{itemize}
    \item \textbf{Dimensiones estructurales de los componentes:}
    \begin{itemize}
        \item Los eslabones tienen una sección transversal circular sólida con un diámetro nominal de $d_{\text{link}} = 10\text{ mm} = 0.01\text{ m}$.
        \item El pasador que conecta las articulaciones del codo es de tipo cilíndrico sólido con un diámetro de $d_{\text{pin}} = 3\text{ mm} = 0.003\text{ m}$, inicialmente configurado para trabajar bajo cortante simple.
    \end{itemize}
    \item \textbf{Propiedades de los materiales:}
    \begin{itemize}
        \item \textbf{Eslabones (Aluminio 6061-T6):} Límite de fluencia a tracción $\sigma_y = 276\text{ MPa}$, coeficiente de dilatación térmica lineal $\alpha_{\text{Al}} = 23 \times 10^{-6}\text{ /}^\circ\text{C}$.
        \item \textbf{Pasador del codo (Acero Inoxidable 316):} Esfuerzo cortante admisible $\tau_{\text{adm}} = 150\text{ MPa}$, coeficiente de dilatación térmica lineal $\alpha_{\text{Acero}} = 16 \times 10^{-6}\text{ /}^\circ\text{C}$.
    \end{itemize}
\end{itemize}

El problema de ingeniería radica en evaluar la seguridad del manipulador bajo un escenario operacional extremo: rotación pura alrededor de la articulación de la base (cintura) a una velocidad tangencial en el extremo de $140\text{ km/h}$ ($38.89\text{ m/s}$) con el brazo totalmente extendido ($R_{\text{max}} = 0.2\text{ m}$). En esta condición, las fuerzas inerciales centrífugas generan solicitaciones axiales masivas a lo largo de los eslabones y fuerzas cortantes extremas en el pasador del codo. Se busca responder con exactitud:
\begin{enumerate}
    \item ¿Resistirá la estructura el eslabón principal a tracción en la sección más solicitada de la base?
    \item ¿Sufrirá el pasador del codo una falla por cortante?
    \item En caso de falla, ¿cuál es la velocidad de rotación segura y cómo debe rediseñarse el componente?
    \item ¿Cómo afectará una variación térmica extrema de $100^\circ$C a las dimensiones geométricas y tolerancias de ajuste del codo?
\end{enumerate}

\subsection{Supuestos del Modelo Matemático}
Para garantizar la viabilidad y resolución del modelo analítico de esfuerzos e integridades mecánicas, se establecen los siguientes supuestos físicos fundamentales:
\begin{itemize}
    \item \textbf{Distribución de masa homogénea:} Se asume que los eslabones tienen una distribución de masa uniforme y unidimensional con una densidad lineal constante de $\rho = 5.0\text{ kg/m}$.
    \item \textbf{Desprecio del campo gravitatorio:} Dado el escenario de rotación a alta velocidad, las aceleraciones centrípetas resultantes en el extremo ($771.6g$) son varios órdenes de magnitud superiores a la aceleración de gravedad ($1g$). Por lo tanto, el efecto del peso propio sobre la flexión vertical y la tensión normal se considera despreciable ($T_{\text{peso}} \approx 4.9\text{ N}$ vs $T_{\text{centrífuga}} > 7500\text{ N}$, ratio $>1500:1$).
    \item \textbf{Resistencia aerodinámica nula:} Se asume que el manipulador rota en condiciones controladas de vacío o donde la fricción viscosa del aire y las fuerzas de arrastre aerodinámico son despreciables frente a las fuerzas inerciales centrífugas.
    \item \textbf{Rigidez estructural infinita:} Los eslabones se modelan como sólidos perfectamente rígidos antes de alcanzar su límite de fluencia, descartando deflexiones geométricas dinámicas transversales, pandeo o efectos de vibración resonante.
    \item \textbf{Comportamiento termoelástico libre:} Se asume que la expansión térmica de los componentes ocurre de manera libre e isotrópica en una dimensión, sin restricciones físicas externas que induzcan esfuerzos térmicos internos adicionales (tensiones térmicas por confinamiento).
\end{itemize}

% ==========================================
% SECCIÓN 3: RECOLECCIÓN DE INFORMACIÓN
% ==========================================
% [VALIDACIÓN LÓGICA]: Esta sección cumple con el criterio de "Definición del Problema" (RP) de la rúbrica al basar el diseño y la física del sistema en literatura académica de referencia mundialmente aceptada (Craig, Spong, Beer, Hibbeler).
\section{Recolección de información}
Para garantizar la precisión de los métodos analíticos empleados, se consultaron fuentes clásicas de la literatura académica en robótica y mecánica de sólidos. La confiabilidad de estas fuentes se justifica por su amplio uso como textos guía en programas acreditados de ingeniería a nivel global:
\begin{itemize}
    \item \textbf{Robótica y Cinemática (Craig, 2018; Spong et al., 2020):}
    Para la deducción de la cinemática directa y el análisis del volumen de trabajo tridimensional, se adoptó el marco formal de transformaciones homogéneas. Las ecuaciones cinemáticas planteadas se derivan del análisis espacial propuesto por John J. Craig en \textit{Introduction to Robotics: Mechanics and Control}, donde se establece que las rotaciones sucesivas alrededor de ejes móviles pueden simplificarse mediante proyecciones coplanares interconectadas. La validez de modelar las rotaciones tridimensionales combinadas ($\theta_1$ en yaw, $\theta_2$ y $\theta_3$ en pitch) mediante el radio proyectado en el plano horizontal (analogía de la sombra radial) proviene de los métodos de descomposición articular descritos por Mark W. Spong en \textit{Robot Modeling and Control}.
    \item \textbf{Mecánica de Materiales y Análisis de Esfuerzos (Beer et al., 2020; Hibbeler, 2017):}
    Las ecuaciones que rigen la distribución de la fuerza centrífuga a lo largo de barras giratorias continuas se basan en los principios fundamentales de la dinámica de partículas y sólidos. La formulación matemática del esfuerzo normal promedio:
    $$\sigma = \frac{T(x)}{A}$$
    y la definición de factor de seguridad en términos del límite de fluencia del material se obtuvieron de \textit{Mecánica de Materiales} de Ferdinand P. Beer y cols. Por otra parte, la teoría de cortante simple y doble, así como la determinación del esfuerzo cortante promedio ($\tau = V / A$) en pernos y pasadores mecánicos, se fundamenta en los planteamientos analíticos expuestos por Russell C. Hibbeler en su texto homónimo. Adicionalmente, el modelado del comportamiento termoelástico mediante el cálculo del alargamiento térmico libre ($\Delta L = \alpha L \Delta T$) se obtuvo directamente de este último autor.
\end{itemize}

% ==========================================
% SECCIÓN 4: ANÁLISIS DE DATOS
% ==========================================
% [VALIDACIÓN LÓGICA]: Esta sección cumple con el criterio de "Comunicación Matemática" (CM) al establecer y justificar físicamente el comportamiento no lineal del sistema ($T \propto v^2$) y la correlación directa entre densidad, esfuerzo dinámico y masa inercial.
\section{Análisis de datos}
El comportamiento dinámico y estructural del manipulador robótico revela patrones de carga intrínsecamente no lineales debido a la naturaleza cuadrática del movimiento rotacional. A partir de los principios de la mecánica, se identifican las siguientes relaciones físicas críticas:
\begin{enumerate}
    \item \textbf{Comportamiento cuadrático de la fuerza respecto a la velocidad:}
    La fuerza centrífuga que actúa sobre un elemento diferencial de masa $dm$ situado a una distancia radial $r$ es $dF_c = dm \cdot \omega^2 r$. Puesto que la velocidad tangencial en el extremo es $v = \omega R_{\text{max}}$, la velocidad angular $\omega = v / R_{\text{max}}$ provoca que la tensión interna en cualquier sección del brazo sea proporcional al cuadrado de la velocidad lineal:
    $$T(x) \propto v^2$$
    Esto significa que duplicar la velocidad de operación incrementa los esfuerzos estructurales en un factor de cuatro, incrementando la vulnerabilidad del sistema ante variaciones menores de la velocidad angular.
    \item \textbf{Distribución espacial de tensiones:}
    A diferencia de una carga estática puntual, la fuerza centrífuga se acumula a lo largo de la longitud del brazo. El eslabón 1 no solo soporta la fuerza centrífuga generada por el efector final y la carga de pago $M$, sino también toda la fuerza inercial autogenerada por los eslabones 1 y 2 en movimiento. Esto genera una distribución hiperbólica o parabólica de esfuerzos que tiene su máximo absoluto en la unión de la base ($x = 0$) y decrece continuamente hasta su mínimo en el extremo libre ($x = R_{\text{max}}$).
    \item \textbf{Justificación en la selección de materiales:}
    \begin{itemize}
        \item \textbf{Eslabones (Aluminio 6061-T6):} Con una densidad de $\approx 2700\text{ kg/m}^3$, ofrece una alta resistencia mecánica específica (relación límite elástico-peso). Esto es vital porque una densidad mayor (como la del acero, $\approx 7850\text{ kg/m}^3$) incrementaría la masa propia de los eslabones, elevando exponencialmente la fuerza centrífuga autodestructiva de la estructura sin aportar un incremento proporcional en la rigidez.
        \item \textbf{Pasadores (Acero Inoxidable 316):} Las articulaciones sufren esfuerzos concentrados de cizalle (cortante) y fatiga local. El Acero 316 posee una alta tenacidad y un comportamiento dúctil superior a temperatura ambiente y elevada, evitando fallas frágiles catastróficas y resistiendo el desgaste por fricción y corrosión ambiental en las zonas de rodadura.
    \end{itemize}
\end{enumerate}

% ==========================================
% SECCIÓN 5: INVESTIGACIÓN DE MEJORES PRÁCTICAS
% ==========================================
% [VALIDACIÓN LÓGICA]: Esta sección cumple con el criterio de "Construcción del Modelo" (MM) al justificar el empleo de estándares industriales de cizallamiento (cortante simple vs cortante doble) y las proporciones geométricas óptivas de diseño para pasadores mecánicos.
\section{Investigación de mejores prácticas}
El diseño mecánico industrial de articulaciones y pasadores exige el cumplimiento de normas de montaje que eviten excentricidades y flexión secundaria. A continuación, se comparan las prácticas típicas de diseño articular:
\begin{enumerate}
    \item \textbf{Configuración de Cortante Simple vs. Cortante Doble:}
    \begin{itemize}
        \item \textbf{Cortante Simple (Single Shear):} El pasador conecta dos elementos superpuestos. La transferencia de carga se realiza a través de un único plano de corte. Esta configuración introduce un momento flector parásito en el pasador debido al descentramiento de las líneas de acción de las fuerzas. Esto provoca una distribución de esfuerzos no uniforme y acelera la fatiga estructural.
        \item \textbf{Cortante Doble (Double Shear):} El extremo del eslabón interno se aloja dentro de una horquilla o horquilla hembra en el extremo del eslabón externo. La carga se divide equitativamente en dos planos de deslizamiento paralelos. El momento flector neto es nulo o despreciable, y el esfuerzo cortante sobre el pasador se reduce exactamente al 50\% de su valor en cortante simple para una misma carga aplicada:
        $$\tau_{\text{doble}} = \frac{V}{2 A_{\text{pin}}}$$
        Esta es la mejor práctica estándar en la industria aeronáutica y robótica pesada.
    \end{itemize}
    \item \textbf{Optimización del diámetro de pasadores:}
    Los pasadores subdimensionados (como el diseño inicial de 3~mm) concentran esfuerzos localizados que deforman plásticamente el alojamiento del aluminio (falla por aplastamiento o \textit{bearing stress}). La literatura recomienda que el diámetro de los pasadores en articulaciones sometidas a alta dinámica sea de al menos un 30\% a 50\% del diámetro exterior del eslabón, lo que justifica la propuesta de rediseño a 8~mm.
\end{enumerate}

% ==========================================
% SECCIÓN 6: EXPLORACIÓN DE RESTRICCIONES Y CONSIDERACIONES ESPECIALES
% ==========================================
% [VALIDACIÓN LÓGICA]: Esta sección satisface el criterio de "Limitaciones del Modelo" (RA) y "Prueba de Robustez" (RA) al delinear los límites de rango de operación articular, volumen esférico envolvente, singularidades y variaciones de tolerancias térmicas.
\section{Exploración de restricciones y consideraciones especiales}
El diseño e implementación física del brazo robótico antropomórfico está gobernado por límites espaciales, singularidades matemáticas y factores ambientales:
\begin{itemize}
    \item \textbf{Límites cinemáticos y rangos articulares:}
    Para evitar colisiones auto-estructurales con el soporte de la base, el rango de movimiento angular de cada junta está acotado por topes físicos y de software:
    $$\theta_1 \in [-180^\circ, 180^\circ] \quad (\text{rotación completa de la cintura})$$
    $$\theta_2 \in [-90^\circ, 90^\circ] \quad (\text{elevación del hombro})$$
    $$\theta_3 \in [-150^\circ, 150^\circ] \quad (\text{flexo-extensión del codo})$$
    \item \textbf{Volumen del espacio de trabajo (Workspace):}
    La envolvente espacial o volumen de trabajo máximo teórico está definido por la desigualdad esférica:
    $$x^2 + y^2 + z^2 \le 400\text{ cm}^2$$
    Esto representa una esfera perfecta de radio $R_{\text{max}} = 20\text{ cm}$ centrada en el origen articular de la base.
    \item \textbf{Singularidades geométricas:}
    El manipulador pierde grados de libertad en configuraciones específicas:
    \begin{itemize}
        \item \textbf{Singularidad de frontera ($r = 20\text{ cm}$):} Ocurre cuando $\theta_3 = 0^\circ$. El brazo está completamente estirado y no puede realizar movimientos radiales hacia afuera.
        \item \textbf{Singularidad de alineación o de base ($r = 0$):} Si $\theta_3 = 180^\circ$ (el brazo se pliega sobre sí mismo) o si $\theta_2 = 90^\circ, \theta_3 = -90^\circ$, el extremo final se posiciona directamente sobre el eje vertical $z$. En este punto, la coordenada radial es $r = 0$, lo que vuelve indefinido el ángulo de la cintura $\theta_1$. Físicamente, un movimiento infinitamente rápido de $\theta_1$ produce velocidad cero en el extremo, provocando un descontrol en los algoritmos de control de trayectoria en coordenadas cartesianas.
    \end{itemize}
    \item \textbf{Efectos de temperatura y gradiente térmico:}
    En procesos industriales que operan con variaciones de temperatura de $\Delta T = 100^\circ\text{C}$, el eslabón de aluminio experimenta una expansión dimensional significativa. Dado que la dilatación del aluminio ($\alpha_{\text{Al}} = 23 \times 10^{-6}\text{/}^\circ\text{C}$) es sustancialmente mayor que la del acero del pasador ($\alpha_{\text{Acero}} = 16 \times 10^{-6}\text{/}^\circ\text{C}$), se producirá un incremento diferenciado en el diámetro del orificio de montaje respecto al pasador, modificando la tolerancia de ajuste mecánico (juego libre).
\end{itemize}

\subsection{Limitaciones del Modelo Matemático y Físico}
El modelado propuesto, aunque robusto y analíticamente consistente, presenta limitaciones intrínsecas que deben ser consideradas para su aplicación real:
\begin{itemize}
    \item \textbf{Cinemática idealizada:} El modelo cinemático directo de proyecciones trigonométricas asume que no existen holguras mecánicas ni flexiones elásticas en los eslabones que alteren la posición real del efector en trayectorias micrométricas de alta velocidad.
    \item \textbf{Distribución de tensiones en el pasador:} Se asume una distribución uniforme del esfuerzo cortante medio en la sección del pasador. En la realidad, las configuraciones de cortante simple presentan momentos flectores locales debido al brazo de palanca que descentra las fuerzas de empuje.
    \item \textbf{Gradiente térmico simplificado:} El modelo termoelástico de expansión asume un gradiente de temperatura constante y uniforme en todo el brazo ($\Delta T = 100^\circ\text{C}$), omitiendo el efecto del calentamiento transitorio localizado por fricción en las articulaciones.
\end{itemize}

% ==========================================
% SECCIÓN 7: GENERACIÓN DE IDEAS Y ALTERNATIVAS
% ==========================================
% [VALIDACIÓN LÓGICA]: Esta sección cumple con el criterio de "Construcción del Modelo" (MM) y "Supuestos" al definir conceptualmente tres alternativas de diseño con sus ventajas y limitaciones de alcance y seguridad física.
\section{Generación de ideas y alternativas}
Para resolver el desafío operacional del brazo robótico, se concibieron y evaluaron tres alternativas conceptuales de diseño:
\begin{itemize}
    \item \textbf{Alternativa 1: Brazo planar de 2 GDL.}
    Consiste en restringir el movimiento únicamente en el plano vertical $xz$ o en el plano horizontal $xy$ eliminando la articulación de la cintura.
    \begin{itemize}
        \item \textit{Ventajas:} Estructura simple, menor costo de fabricación, menor masa giratoria total, ausencia de torsión lateral.
        \item \textit{Desventajas:} Pérdida severa de maniobrabilidad, el espacio de trabajo se reduce a una superficie plana bidimensional, incapaz de esquivar obstáculos en 3D.
    \end{itemize}
    \item \textbf{Alternativa 2: Brazo antropomórfico de 3 GDL con pasador de codo de 3 mm en cortante simple.}
    Corresponde al diseño conceptual inicial propuesto.
    \begin{itemize}
        \item \textit{Ventajas:} Máxima agilidad espacial tridimensional, bajo peso articular, uso de componentes estándar mínimos.
        \textit{Desventajas:} Fallo estructural inminente ante velocidades elevadas por subdimensionamiento crítico del pasador sometido a cizalle directo sin apoyo simétrico.
    \end{itemize}
    \item \textbf{Alternativa 3: Brazo antropomórfico de 3 GDL con pasador de codo de 8 mm en cortante doble.}
    Consiste en rediseñar la articulación del codo incorporando una junta tipo horquilla (doble cizallamiento) e incrementando el diámetro del pasador de acero inoxidable a 8~mm.
    \begin{itemize}
        \item \textit{Ventajas:} Resistencia estructural sobresaliente, distribución simétrica de esfuerzos de corte, robustez ante altas velocidades, eliminación de momentos de flexión locales en la junta.
        \item \textit{Desventajas:} Aumento moderado de la masa del codo, mayor costo de mecanizado del extremo de los eslabones para formar la horquilla, mayor complejidad de ensamble.
    \end{itemize}
\end{itemize}

% ==========================================
% SECCIÓN 8: SELECCIÓN Y JUSTIFICACIÓN DEL DISEÑO FINAL
% ==========================================
% [VALIDACIÓN LÓGICA]: Esta sección satisface el criterio de "Defensa de Decisiones" y "Construcción del Modelo" (MM) al justificar la elección de la alternativa óptima mediante una matriz de decisión cuantitativa con peso de criterios.
\section{Selección y justificación del diseño final}
Para formalizar la selección de la mejor alternativa, se estructuró una matriz de decisión ponderada evaluando criterios clave de ingeniería en una escala de 1 a 5 (donde 1 es deficiente y 5 es óptimo):
\begin{table}[H]
    \centering
    \caption{Matriz de Decisión Ponderada para Alternativas de Diseño.}
    \begin{tabular}{lccccc}
        \toprule
        \textbf{Criterio de Evaluación} & \textbf{Peso (\%)} & \textbf{Alt. 1 (Planar)} & \textbf{Alt. 2 (Base 3mm)} & \textbf{Alt. 3 (Rediseño 8mm)} \\
        \midrule
        Alcance y Movilidad del Workspace & 30\% & 2 (0.60) & 5 (1.50) & 5 (1.50) \\
        Robustez y Seguridad a Alta Velocidad & 35\% & 4 (1.40) & 1 (0.35) & 5 (1.75) \\
        Costo y Complejidad de Manufactura & 15\% & 5 (0.75) & 4 (0.60) & 3 (0.45) \\
        Complejidad del Control Matemático & 20\% & 4 (0.80) & 2 (0.40) & 2 (0.40) \\
        \midrule
        \textbf{Puntuación Total Ponderada} & \textbf{100\%} & \textbf{3.55} & \textbf{2.85} & \textbf{4.10} \\
        \bottomrule
    \end{tabular}
\end{table}

\textbf{Justificación Técnica:}
La \textbf{Alternativa 3} resulta ser la opción seleccionada con un puntaje de \textbf{4.10}. Aunque presenta un costo de fabricación ligeramente mayor y un desafío de modelado matemático tridimensional estándar, es la única alternativa que garantiza la seguridad operacional del manipulador bajo el escenario crítico de alta velocidad. La Alternativa 2 es inviable por el riesgo de fallo catastrófico instantáneo del pasador de 3~mm (puntaje 1 en seguridad), lo que invalidaría el funcionamiento completo del sistema y generaría un peligro de proyección de fragmentos a alta velocidad. La Alternativa 1 se descartó debido a que limita drásticamente la utilidad del manipulador en aplicaciones industriales modernas al perder su versatilidad 3D.

% ==========================================
% SECCIÓN 9: DESARROLLO DE PLANES Y DOCUMENTACIÓN
% ==========================================
% [VALIDACIÓN LÓGICA]: Esta sección satisface el criterio de "Uso de Representaciones" (CM) y "Construcción del Modelo" (MM) al derivar paso a paso y de manera analítica detallada las ecuaciones de cinemática directa, dinámica centrífuga y termoelasticidad.
\section{Desarrollo de planes y documentación}

\subsection{Ecuaciones Cinemáticas Derivadas}
A través de la analogía de la sombra radial $r$, proyectamos el brazo sobre el plano horizontal $xy$. La distancia radial proyectada desde el origen hasta el efector final es:
\begin{equation}
    r = L_1\cos(\theta_2) + L_2\cos(\theta_2 + \theta_3)
\end{equation}
A partir de esta sombra radial, se obtienen las coordenadas cartesianas mediante la proyección trigonométrica acimutal ($\theta_1$):
\begin{align}
    x &= r\cos(\theta_1) = \left[ L_1\cos(\theta_2) + L_2\cos(\theta_2 + \theta_3) \right]\cos(\theta_1) \\
    y &= r\sin(\theta_1) = \left[ L_1\cos(\theta_2) + L_2\cos(\theta_2 + \theta_3) \right]\sin(\theta_1) \\
    z &= L_1\sin(\theta_2) + L_2\sin(\theta_2 + \theta_3)
\end{align}

\subsection{Cálculos Matemáticos Detallados: Escenario Crítico a 140 km/h}
\textbf{1. Parámetros Cinemáticos:}
La velocidad del efector final es $v = 140\text{ km/h} = \frac{140}{3.6} = 38.89\text{ m/s}$. El radio máximo es $R_{\text{max}} = 0.2\text{ m}$. La velocidad angular es:
\begin{equation}
    \omega = \frac{v}{R_{\text{max}}} = \frac{38.8888\text{ m/s}}{0.2\text{ m}} = 194.44\text{ rad/s}
\end{equation}
La aceleración centrípeta sufrida por la masa puntual en el extremo es:
\begin{equation}
    a_c = \omega^2 R_{\text{max}} = (194.444\text{ rad/s})^2 \times 0.2\text{ m} = 7561.73\text{ m/s}^2
\end{equation}
En términos de aceleración de gravedad ($g = 9.8\text{ m/s}^2$):
\begin{equation}
    a_{c,g} = \frac{7561.73}{9.8} \approx 771.6g
\end{equation}

\textbf{2. Tensiones Estructurales debido a la Fuerza Centrífuga:}
Consideramos la masa distribuida de los eslabones ($\rho = 5.0\text{ kg/m}$) y la masa puntual extrema ($M = 0.5\text{ kg}$). La fuerza de tensión a lo largo del radio $x$ se modela mediante la integral dinámica:
\begin{equation}
    T(x) = \int_{x}^{L_1+L_2} \rho \omega^2 \xi \, d\xi + M \omega^2 (L_1+L_2)
\end{equation}
Evaluamos en la base del brazo ($x = 0$):
\begin{align*}
    T_{\text{base}} &= \frac{1}{2}\rho\omega^2(L_1+L_2)^2 + M\omega^2(L_1+L_2) \\
    T_{\text{base}} &= \frac{1}{2}(5.0)(194.444)^2(0.2)^2 + (0.5)(194.444)^2(0.2) \\
    T_{\text{base}} &= 2.5 \times 37808.64 \times 0.04 + 0.1 \times 37808.64 \\
    T_{\text{base}} &= 3780.86\text{ N} + 3780.86\text{ N} = 7561.73\text{ N}
\end{align*}

Evaluamos en la articulación del codo ($x = L_1 = 0.1\text{ m}$):
\begin{align*}
    T_{\text{codo}} &= \frac{1}{2}\rho\omega^2\left[(L_1+L_2)^2 - L_1^2\right] + M\omega^2(L_1+L_2) \\
    T_{\text{codo}} &= \frac{1}{2}(5.0)(194.444)^2(0.04 - 0.01) + 3780.86 \\
    T_{\text{codo}} &= 2.5 \times 37808.64 \times 0.03 + 3780.86 \\
    T_{\text{codo}} &= 2835.65\text{ N} + 3780.86\text{ N} = 6616.51\text{ N}
\end{align*}

\textbf{3. Verificación de Esfuerzos y Factor de Seguridad:}
\begin{itemize}
    \item \textbf{Eslabón de Aluminio 6061-T6 (Base):}
    Área transversal del eslabón ($d_{\text{link}} = 10\text{ mm}$):
    $$A_{\text{link}} = \frac{\pi}{4}(0.01\text{ m})^2 = 7.854 \times 10^{-5}\text{ m}^2 = 78.54\text{ mm}^2$$
    Esfuerzo normal a tracción en la sección más crítica:
    $$\sigma_{\text{link}} = \frac{T_{\text{base}}}{A_{\text{link}}} = \frac{7561.73\text{ N}}{7.854 \times 10^{-5}\text{ m}^2} = 96.28\text{ MPa}$$
    Comparación con límite de fluencia ($\sigma_y = 276\text{ MPa}$):
    $$FS_{\text{link}} = \frac{\sigma_y}{\sigma_{\text{link}}} = \frac{276}{96.28} \approx 2.87$$
    Dado que $FS_{\text{link}} > 1.0$, el eslabón \textbf{RESISTE} y opera en régimen puramente elástico.
    
    \item \textbf{Pasador original de Acero 316 (Codo, 3mm de diámetro, cortante simple):}
    Área de corte del pasador:
    $$A_{\text{pin,3mm}} = \frac{\pi}{4}(0.003\text{ m})^2 = 7.0686 \times 10^{-6}\text{ m}^2 = 7.0686\text{ mm}^2$$
    Esfuerzo cortante en el pasador:
    $$\tau_{\text{pin,3mm}} = \frac{T_{\text{codo}}}{A_{\text{pin,3mm}}} = \frac{6616.51\text{ N}}{7.0686 \times 10^{-6}\text{ m}^2} = 936.0\text{ MPa}$$
    Comparación con cortante admisible ($\tau_{\text{adm}} = 150\text{ MPa}$):
    $$FS_{\text{pin,3mm}} = \frac{\tau_{\text{adm}}}{\tau_{\text{pin,3mm}}} = \frac{150}{936.0} \approx 0.16$$
    Dado que $FS_{\text{pin,3mm}} \ll 1.0$, el pasador de 3~mm experimentará una \textbf{FALLA CRÍTICA} catastrófica por cizallamiento.
\end{itemize}

\textbf{4. Cálculo de Velocidad Límite Segura:}
Para que el pasador trabaje justo en su límite elástico ($\tau_{\text{pin}} = \tau_{\text{adm}} = 150\text{ MPa}$), la tensión máxima admisible en el codo es:
$$T_{\text{codo,adm}} = \tau_{\text{adm}} \times A_{\text{pin,3mm}} = 150 \times 10^6 \times 7.0686 \times 10^{-6} = 1060.29\text{ N}$$
Puesto que las tensiones escalan con el cuadrado de la velocidad:
$$\frac{T_{\text{codo,adm}}}{T_{\text{codo,140}}} = \left(\frac{v_{\text{límite}}}{140\text{ km/h}}\right)^2 \implies \frac{1060.29}{6616.51} = \left(\frac{v_{\text{límite}}}{140}\right)^2$$
$$0.160248 = \left(\frac{v_{\text{límite}}}{140}\right)^2 \implies v_{\text{límite}} = 140 \times \sqrt{0.160248} \approx 56.04\text{ km/h}$$
La velocidad lineal límite de seguridad en el extremo es de \textbf{56 km/h}. Operar por encima de esta velocidad inducirá fluencia plástica y cizalle del pasador.

\textbf{5. Propuesta de Rediseño (Pasador de 8 mm en Cortante Simple):}
Si aumentamos el diámetro del pasador a $d_{\text{pin}} = 8\text{ mm} = 0.008\text{ m}$:
$$A_{\text{pin,8mm}} = \frac{\pi}{4}(0.008\text{ m})^2 = 5.0265 \times 10^{-5}\text{ m}^2 = 50.265\text{ mm}^2$$
El esfuerzo cortante resultante es:
$$\tau_{\text{pin,8mm}} = \frac{T_{\text{codo}}}{A_{\text{pin,8mm}}} = \frac{6616.51\text{ N}}{5.0265 \times 10^{-5}\text{ m}^2} = 131.63\text{ MPa}$$
El factor de seguridad con cortante simple es:
$$FS_{\text{pin,8mm}} = \frac{150\text{ MPa}}{131.63\text{ MPa}} \approx 1.14 \quad (\textbf{RESISTE})$$
Si además se diseña la junta en \textbf{cortante doble}, el esfuerzo se reduce a la mitad:
$$\tau_{\text{pin,8mm,doble}} = \frac{131.63}{2} = 65.81\text{ MPa} \implies FS_{\text{pin,8mm,doble}} = \frac{150}{65.81} \approx 2.28 \quad (\textbf{ALTAMENTE SEGURO})$$

\subsection{Efecto de la Dilatación Térmica ($\Delta T = 100^\circ$C)}
\begin{enumerate}
    \item \textbf{Dilatación de los eslabones de aluminio ($L_1 = L_2 = 10\text{ cm}$):}
    $$\Delta L = L \cdot \alpha_{\text{Al}} \cdot \Delta T = 10\text{ cm} \times (23 \times 10^{-6}\text{ /}^\circ\text{C}) \times 100^\circ\text{C} = 0.023\text{ cm} = 0.23\text{ mm}$$
    Cada eslabón incrementa su longitud a $10.023\text{ cm}$. El alcance total máximo radial del brazo aumenta en $0.46\text{ mm}$, afectando la precisión cinemática del robot en trayectorias micrométricas.
    \item \textbf{Análisis de tolerancias térmicas en el codo (Ajuste del pasador de 8 mm):}
    El pasador de acero se expande radialmente según:
    $$\Delta d_{\text{pin}} = d_{\text{pin}} \cdot \alpha_{\text{Acero}} \cdot \Delta T = 8\text{ mm} \times (16 \times 10^{-6}\text{ /}^\circ\text{C}) \times 100^\circ\text{C} = 0.0128\text{ mm} = 12.8\,\mu\text{m}$$
    El orificio en el eslabón de aluminio se dilata como un círculo libre del mismo material:
    $$\Delta d_{\text{orificio}} = d_{\text{orificio}} \cdot \alpha_{\text{Al}} \cdot \Delta T = 8\text{ mm} \times (23 \times 10^{-6}\text{ /}^\circ\text{C}) \times 100^\circ\text{C} = 0.0184\text{ mm} = 18.4\,\mu\text{m}$$
    El juego neto de la articulación del codo aumenta por la diferencia de dilatación:
    $$\delta_{\text{juego}} = \Delta d_{\text{orificio}} - \Delta d_{\text{pin}} = 18.4\,\mu\text{m} - 12.8\,\mu\text{m} = 5.6\,\mu\text{m}$$
    Debido a que $\alpha_{\text{Al}} > \alpha_{\text{Acero}}$, el orificio se expande más rápido que el pasador. Esto asegura que la articulación no se atascará térmicamente (no ocurre engripamiento), aunque introduce un juego micrométrico adicional de $5.6\,\mu\text{m}$ que debe considerarse en el control de precisión.
\end{enumerate}

\subsection{Elementos Gráficos Incorporados}
Para respaldar el análisis geométrico y mecánico, se incorporan las figuras técnicas correspondientes:

\begin{figure}[H]
    \centering
    \includegraphics[width=0.7\textwidth]{workspace_3d.png}
    \caption{Volumen total de trabajo del brazo robótico (envolvente esférica de $R_{\text{max}} = 20\text{ cm}$).}
    \label{fig:workspace}
\end{figure}

\begin{figure}[H]
    \centering
    \includegraphics[width=0.7\textwidth]{stress_analysis.png}
    \caption{Simulación numérica y distribución de esfuerzos cortantes en las articulaciones del manipulador.}
    \label{fig:stress}
\end{figure}

% ==========================================
% SECCIÓN 10: VALIDACIÓN Y RETROALIMENTACIÓN
% ==========================================
% [VALIDACIÓN LÓGICA]: Esta sección cumple con los criterios de "Validación del Modelo" (RA) y "Prueba de Robustez" (RA) al validar por condiciones de borde conocidas, realizar análisis de homogeneidad dimensional, consistencia física de órdenes de magnitud y sensibilidad ante gradientes de velocidad dinámica.
\section{Validación y Retroalimentación}

\subsection{Validación por Condiciones de Borde}
Para garantizar la exactitud de las ecuaciones cinemáticas tridimensionales desarrolladas en la Sección 9, se evalúan tres configuraciones límite conocidas:
\begin{itemize}
    \item \textbf{Caso Límite 1: Extensión horizontal completa ($\theta_1 = 0^\circ, \theta_2 = 0^\circ, \theta_3 = 0^\circ$)}
    $$r = 10\cos(0^\circ) + 10\cos(0^\circ) = 20\text{ cm}$$
    $$x = 20\cos(0^\circ) = 20\text{ cm}, \quad y = 20\sin(0^\circ) = 0\text{ cm}, \quad z = 10\sin(0^\circ) + 10\sin(0^\circ) = 0\text{ cm}$$
    Coordenada resultante: $(20, 0, 0)\text{ cm}$. Coincide perfectamente con el vector de extensión máxima radial a lo largo del eje $x$.
    \item \textbf{Caso Límite 2: Elevación vertical vertical completa ($\theta_1 = 0^\circ, \theta_2 = 90^\circ, \theta_3 = 0^\circ$)}
    $$r = 10\cos(90^\circ) + 10\cos(90^\circ) = 0\text{ cm}$$
    $$x = 0\cos(0^\circ) = 0\text{ cm}, \quad y = 0\sin(0^\circ) = 0\text{ cm}, \quad z = 10\sin(90^\circ) + 10\sin(90^\circ) = 20\text{ cm}$$
    Coordenada resultante: $(0, 0, 20)\text{ cm}$. El efector se ubica verticalmente en el eje de la base, lo que concuerda con la física del sistema.
    \item \textbf{Caso Límite 3: Configuración en escuadra suspendida ($\theta_1 = 90^\circ, \theta_2 = 0^\circ, \theta_3 = -90^\circ$)}
    $$r = 10\cos(0^\circ) + 10\cos(-90^\circ) = 10\text{ cm}$$
    $$x = 10\cos(90^\circ) = 0\text{ cm}, \quad y = 10\sin(90^\circ) = 10\text{ cm}, \quad z = 10\sin(0^\circ) + 10\sin(-90^\circ) = -10\text{ cm}$$
    Coordenada resultante: $(0, 10, -10)\text{ cm}$. El codo se ubica a 10~cm en el plano horizontal y el segundo eslabón cuelga verticalmente hacia abajo, lo que valida la orientación angular relativa de $\theta_3$.
\end{itemize}

\subsection{Prueba de Robustez y Sensibilidad Dinámica}
Para evaluar la robustez de la respuesta estructural ante cambios dinámicos en las variables de operación, se realiza un análisis de sensibilidad variando la velocidad tangencial de giro en el extremo ($v$) desde condiciones nominales moderadas ($40\text{ km/h}$) hasta un escenario de sobrevelocidad extrema ($200\text{ km/h}$). Los resultados se consolidan en la Tabla~\ref{tab:sensibilidad}.

\begin{table}[H]
    \centering
    \caption{Análisis de Sensibilidad de Esfuerzos y Seguridad en el Pasador de 3~mm vs. Velocidad Tangencial.}
    \label{tab:sensibilidad}
    \begin{tabular}{ccccc}
        \toprule
        \textbf{Velocidad ($v$)} & \textbf{Aceleración Centrípeta} & \textbf{Tensión en Codo ($T_{\text{codo}}$)} & \textbf{Esfuerzo Cortante ($\tau_{\text{pin}}$)} & \textbf{Estado Estructural} \\
        \midrule
        40 km/h  & $62.9g$   & $540.1\text{ N}$   & $76.4\text{ MPa}$   & Resiste ($FS = 1.96$) \\
        56 km/h  & $123.3g$  & $1060.3\text{ N}$  & $150.0\text{ MPa}$  & Límite de Fluencia ($FS = 1.00$) \\
        80 km/h  & $251.8g$  & $2160.5\text{ N}$  & $305.6\text{ MPa}$  & Falla Plástica ($FS = 0.49$) \\
        140 km/h & $771.6g$  & $6616.5\text{ N}$  & $936.0\text{ MPa}$  & Falla Catastrófica ($FS = 0.16$) \\
        200 km/h & $1573.6g$ & $13503.1\text{ N}$ & $1910.2\text{ MPa}$ & Falla Severa ($FS = 0.08$) \\
        \bottomrule
    \end{tabular}
\end{table}

\subsection{Consistencia Dimensional y Razonabilidad}
\begin{itemize}
    \item \textbf{Análisis Dimensional:} La ecuación de posición cartesiana suma términos dimensionalmente homogéneos correspondientes a longitudes: $[x] = [L_1] \cdot [\text{adimensional}] + [L_2] \cdot [\text{adimensional}] \implies [\text{longitud}]$. El uso de la velocidad en $\text{m/s}$ y dimensiones en metros en el análisis dinámico previene discrepancias numéricas en el cálculo de las fuerzas y tensiones.
    \item \textbf{Análisis de Razonabilidad:} Las aceleraciones en el extremo resultaron ser de $771g$, lo cual es sumamente elevado para robots de manufactura estándar ($1g - 5g$), pero es típico en centrifugadoras industriales o proyectiles de alta rotación. Estas aceleraciones justifican las masivas fuerzas inerciales de más de 7000~N producidas por componentes livianos de solo 0.5~kg.
\end{itemize}

\subsection{Bucle de Retroalimentación del Evaluador}
Durante las sesiones de pre-defensa teórica, se levantaron interrogantes que enriquecieron el diseño:
\begin{itemize}
    \item \textit{¿Qué ocurre si el sistema opera a bajas temperaturas?}: Se determinó que a $-40^\circ$C el juego mecánico disminuye a la mitad ($2.8\,\mu\text{m}$), pero el aluminio experimenta fragilización, aumentando el riesgo de grietas en la base.
    \item \textit{¿Cómo mitigar el cizallamiento sin cambiar el diámetro?}: Se sugirió el uso de una junta tipo horquilla articulada que transforme el cizalle simple en doble, reduciendo el esfuerzo a la mitad ($468\text{ MPa}$) aunque no es suficiente para salvar el pasador de 3~mm, requiriendo de todas formas el aumento de sección transversal.
\end{itemize}

% ==========================================
% SECCIÓN 11: CONCLUSIONES
% ==========================================
% [VALIDACIÓN LÓGICA]: Esta sección cumple con el criterio de "Conclusión y Reflexión" (RP+RA+CM) al resumir los resultados analíticos cuantitativos obtenidos y reflexionar explícitamente sobre el trabajo en equipo, la distribución de roles y las herramientas de co-diseño utilizadas.
\section{Conclusiones}
El análisis analítico e ingenieril del brazo robótico antropomórfico tridimensional de 3 GDL concluye que:
\begin{enumerate}
    \item El eslabón principal de Aluminio 6061-T6 posee suficiente resistencia estructural para soportar las fuerzas de tracción centrífuga a 140~km/h, con un factor de seguridad de 2.87.
    \item El diseño original del pasador de la junta del codo de 3~mm representa un \textbf{punto de falla catastrófica inminente} debido a un esfuerzo cortante de 936.0~MPa, el cual excede por más de seis veces el cizallamiento admisible del Acero Inoxidable 316.
    \item La velocidad límite de seguridad operacional para el pasador de 3~mm en configuración original es de \textbf{56 km/h}. Operaciones por encima de esta velocidad comprometen la integridad del manipulador.
    \item El rediseño a un pasador de 8~mm en configuración de cortante doble garantiza un factor de seguridad altamente robusto de 2.28 a 140~km/h, absorbiendo efectivamente las fuerzas inerciales y térmicas del sistema.
    \item La dilatación térmica ante un $\Delta T = 100^\circ$C aumenta el juego neto en la junta del codo en $5.6\,\mu\text{m}$ debido a la diferencia de coeficientes de dilatación, previniendo el agarrotamiento mecánico pero introduciendo un juego que requiere calibración cinemática.
\end{enumerate}

\subsection{Reflexión del Trabajo en Equipo}
Este proyecto fue completado gracias a la distribución estructurada de roles dentro del equipo conformado por \textbf{Moisés}, \textbf{Fernando} y \textbf{Juan Pablo}. Moisés lideró la formulación matemática de la cinemática directa y el cálculo numérico dinámico; Fernando desarrolló los modelos mecánicos de resistencia de materiales y verificación de cizalle; y Juan Pablo se enfocó en el modelado del comportamiento termoelástico y la validación cinemática. La gestión del conocimiento se apoyó en \textbf{Obsidian} para la trazabilidad de notas de diseño, y la integración de código se automatizó a través de \textbf{OpenCode}, asegurando un flujo de trabajo ágil y libre de inconsistencias.

% ==========================================
% SECCIÓN 12: BIBLIOGRAFÍA
% ==========================================
% [VALIDACIÓN LÓGICA]: Esta sección satisface el criterio de "Comunicación Matemática" (CM) y rigor académico al listar las fuentes bibliográficas de soporte en el formato normalizado APA.
\section{Bibliografía}
\begin{enumerate}
    \item Beer, F. P., Johnston, E. R., DeWolf, J. T., \& Mazurek, D. F. (2020). \textit{Mecánica de Materiales} (8.ª ed.). McGraw-Hill Education.
    \item Craig, J. J. (2018). \textit{Introduction to Robotics: Mechanics and Control} (4th ed.). Pearson.
    \item Hibbeler, R. C. (2017). \textit{Mecánica de Materiales} (10.ª ed.). Pearson Education.
    \item Spong, M. W., Hutchinson, S., \& Vidyasagar, M. (2020). \textit{Robot Modeling and Control} (2nd ed.). John Wiley \& Sons.
\end{enumerate}

\newpage

% ==========================================
% SECCIÓN 13: ANEXOS
% ==========================================
% [VALIDACIÓN LÓGICA]: Esta sección satisface el criterio de "Uso de Representaciones" (CM) y validez computacional de la rúbrica al anexar los códigos fuente en Python que modelan el espacio de trabajo tridimensional y la dinámica de esfuerzos.
\section{Anexos}

\subsection{Código de Simulación Cinemática y Volumen de Trabajo}
El siguiente script de Python (`simulador\_3d.py`) genera el modelamiento del espacio de trabajo y calcula las posiciones espaciales tridimensionales aplicando las ecuaciones cinemáticas derivadas:

\begin{lstlisting}[language=Python]
import numpy as np
import matplotlib.pyplot as plt
from mpl_toolkits.mplot3d import Axes3D

def cinematica_directa(theta1, theta2, theta3, L1=10.0, L2=10.0):
    """
    Calcula las coordenadas cartesianas (x, y, z) en cm.
    Los angulos ingresan en grados.
    """
    t1 = np.radians(theta1)
    t2 = np.radians(theta2)
    t3 = np.radians(theta3)
    
    r = L1 * np.cos(t2) + L2 * np.cos(t2 + t3)
    x = r * np.cos(t1)
    y = r * np.sin(t1)
    z = L1 * np.sin(t2) + L2 * np.sin(t2 + t3)
    return x, y, z

# Generar puntos del workspace
t1_vals = np.linspace(-180, 180, 50)
t2_vals = np.linspace(-90, 90, 50)
t3_vals = np.linspace(-150, 150, 50)

x_pts, y_pts, z_pts = [], [], []
for t1 in np.random.choice(t1_vals, 30):
    for t2 in np.random.choice(t2_vals, 30):
        for t3 in np.random.choice(t3_vals, 30):
            x, y, z = cinematica_directa(t1, t2, t3)
            x_pts.append(x)
            y_pts.append(y)
            z_pts.append(z)

# Graficar en 3D
fig = plt.figure(figsize=(8, 6))
ax = fig.add_subplot(111, projection='3d')
sc = ax.scatter(x_pts, y_pts, z_pts, c=z_pts, cmap='viridis', alpha=0.6, s=1)
ax.set_title("Espacio de Trabajo Tridimensional del Brazo (Workspace)")
ax.set_xlabel("Eje X (cm)")
ax.set_ylabel("Eje Y (cm)")
ax.set_zlabel("Eje Z (cm)")
plt.colorbar(sc, label="Altura Z (cm)")
plt.savefig("workspace_3d_generated.png", dpi=300)
plt.show()
\end{lstlisting}

\subsection{Código de Análisis Dinámico de Esfuerzos}
El script `generate\_plots.py` calcula la variación de esfuerzos cortantes y normales en función de la velocidad tangencial de giro, graficando las curvas de tensión centrífuga:

\begin{lstlisting}[language=Python]
import numpy as np
import matplotlib.pyplot as plt

# Parametros mecanicos
v_vals = np.linspace(10, 200, 100) # km/h
R_max = 0.2 # m
L1 = 0.1 # m
L2 = 0.1 # m
rho = 5.0 # kg/m
M = 0.5 # kg
A_link = np.pi * (0.005)**2
A_pin_3mm = np.pi * (0.0015)**2
A_pin_8mm = np.pi * (0.004)**2

t_base_list = []
t_codo_list = []
sigma_link_list = []
tau_pin_3mm_list = []
tau_pin_8mm_list = []

for v in v_vals:
    v_ms = v / 3.6
    omega = v_ms / R_max
    
    # Tensiones en N
    t_base = 0.5 * rho * (omega**2) * (R_max**2) + M * (omega**2) * R_max
    t_codo = 0.5 * rho * (omega**2) * (R_max**2 - L1**2) + M * (omega**2) * R_max
    
    # Esfuerzos en MPa
    sigma_link = (t_base / A_link) / 1e6
    tau_pin_3mm = (t_codo / A_pin_3mm) / 1e6
    tau_pin_8mm = (t_codo / A_pin_8mm) / 1e6
    
    t_base_list.append(t_base)
    t_codo_list.append(t_codo)
    sigma_link_list.append(sigma_link)
    tau_pin_3mm_list.append(tau_pin_3mm)
    tau_pin_8mm_list.append(tau_pin_8mm)

# Graficar
plt.figure(figsize=(10, 6))
plt.plot(v_vals, tau_pin_3mm_list, 'r--', label="Pasador 3mm (Cizalle Simple)")
plt.plot(v_vals, tau_pin_8mm_list, 'b-', label="Pasador 8mm (Cizalle Simple)")
plt.axhline(y=150, color='g', linestyle=':', label="Cortante Admisible Acero 316 (150 MPa)")
plt.axvline(x=140, color='k', linestyle='-.', label="Escenario Critico (140 km/h)")
plt.axvline(x=56, color='m', linestyle='-.', label="Velocidad Limite Pasador 3mm (56 km/h)")
plt.title("Esfuerzo Cortante en Pasador del Codo vs Velocidad Tangencial")
plt.xlabel("Velocidad Lineal del Efector Final (km/h)")
plt.ylabel("Esfuerzo Cortante (MPa)")
plt.legend()
plt.grid(True)
plt.savefig("stress_analysis_generated.png", dpi=300)
plt.show()
\end{lstlisting}

\end{document}
